\subsection{Expected Calibration Error}
\label{sec:methods:ece}

A classifier is said to be \textit{calibrated} if its predicted 
confidence matches its empirical accuracy. Formally, for a classifier 
outputting confidence $\hat{p}$ and true label $y$:

\begin{equation}
P(y = \hat{y} \mid \hat{p} = p) = p \quad \forall\, p \in [0, 1]
\label{eq:calibration}
\end{equation}

where $\hat{y}$ is the predicted class. We measure calibration using 
Expected Calibration Error \citep[ECE;][]{Naeini2015}, which 
approximates the expectation in Equation~\ref{eq:calibration} by 
partitioning predictions into $M$ equally spaced confidence bins:

\begin{equation}
\mathrm{ECE} = \sum_{m=1}^{M} \frac{|B_m|}{n} 
\left| \mathrm{acc}(B_m) - \mathrm{conf}(B_m) \right|
\label{eq:ece}
\end{equation}

where $B_m$ is the set of samples whose top-1 confidence falls in bin 
$m$, $n$ is the total number of samples, $\mathrm{acc}(B_m)$ is the 
fraction of correct predictions in bin $m$, and $\mathrm{conf}(B_m)$ 
is the mean confidence in bin $m$. We use $M = 10$ bins throughout. 
A perfectly calibrated classifier has $\mathrm{ECE} = 0$; following
common practice we use 0.05 as a practical reference level for
consequential miscalibration \citep{Guo2017}, while noting this is a
convention, not a formal significance boundary.

For ALeRCE's 15-class classifier, we compute ECE using the top-1
predicted class and its associated probability across all 15 output
classes. For Fink's binary classifiers, which output a scalar
$P(\mathrm{SNIa})$, we apply a binary variant of
Equation~\ref{eq:ece} where $\mathrm{acc}(B_m)$ is the observed
SNIa fraction within bin $m$.

\subsubsection{Per-class calibration: three definitions}
\label{sec:methods:perclass}

``Per-class calibration'' is ambiguous in the multiclass setting, and
the choice of definition materially affects conclusions
(Section~\ref{sec:results:alerce_perclass}). We therefore compute and
report three distinct quantities explicitly:

\begin{enumerate}
\item \textbf{True-class-conditioned top-label ECE}: for class $k$,
Equation~\ref{eq:ece} restricted to objects whose \emph{spectroscopic}
label is $k$. This measures a class-conditional confidence deficit —
how appropriately the classifier weights objects that truly are class
$k$ — which is the operationally relevant quantity when asking whether
genuine members of a class will clear a confidence-based follow-up
threshold. It is \emph{not} calibration in the standard conditional
sense.
\item \textbf{Predicted-class-stratified top-label ECE}:
Equation~\ref{eq:ece} restricted to objects whose \emph{predicted}
top-1 class is $k$. This asks whether confidence matches precision
within each predicted-class stream.
\item \textbf{One-vs-rest classwise ECE}: over \emph{all} objects,
binning the class-$k$ probability $p_k$ against the empirical
frequency of true class $k$ — the standard classwise calibration
condition $P(y = k \mid \hat{p}_k = p) = p$.
\end{enumerate}

\subsubsection{Ranking metrics}
\label{sec:methods:ranking}

ECE alone cannot distinguish a classifier that ranks objects usefully
but reports distorted probabilities (repairable post hoc) from one
with no discriminative signal (not repairable by any monotone
transformation, since monotone recalibration preserves ranking). For
the binary Fink classifiers we therefore also report ROC--AUC and
average precision (PR--AUC), alongside Brier score. A score with
$\mathrm{AUC} \approx 0.5$ can always be made to \emph{appear}
calibrated — a constant prediction at the base rate has
$\mathrm{ECE} \approx 0$ — while carrying no information; reporting
ECE together with ranking metrics guards against this failure mode.

\subsubsection{Robustness to bin construction}
\label{sec:methods:ace}

Equation~\ref{eq:ece}'s equal-width bins are known to be sensitive
to how probability mass is distributed across the confidence range
\citep{Nixon2019, Roelofs2022}: a classifier whose scores concentrate
near one end of $[0,1]$ — Fink's \texttt{rf\_snia\_vs\_nonia} has mean
score $0.011$, putting roughly 90\% of objects in the first of 10
bins — leaves the remaining bins populated by very few points, so the
resulting ECE can be more sensitive to noise or bin placement than a
well-populated case. As a robustness check, we additionally compute
each headline ECE using equal-\textit{mass} (adaptive) bins, i.e.
quantile bins of the confidence/score distribution that guarantee
$\approx n/10$ samples per bin regardless of skew. If equal-width and
equal-mass binning agree, the reported miscalibration is a property
of the classifier and not an artefact of bin construction.

\subsection{Uncertainty Quantification via Bootstrap}
\label{sec:methods:bootstrap}

Point estimates of ECE carry substantial sampling uncertainty, 
particularly for small classes such as SLSN ($n=55$) and TDE ($n=30$). 
We report 95\% confidence intervals on all ECE estimates using the 
nonparametric bootstrap \citep{Efron1979}: we resample the evaluation 
set with replacement $B = 2000$ times, compute ECE on each resample, 
and report the 2.5th and 97.5th percentiles of the resulting 
distribution. All bootstrap procedures use a fixed random seed 
(seed $= 42$) for reproducibility.

\subsection{Reliability Diagrams}
\label{sec:methods:reliability}

We visualise calibration using reliability diagrams \citep{DeGroot1983}, 
which plot observed accuracy against mean predicted confidence for each 
bin. A perfectly calibrated classifier produces points lying on the 
diagonal. Points above the diagonal indicate underconfidence 
(predicted probability lower than empirical accuracy); points below 
indicate overconfidence. We include sample-count histograms alongside 
each diagram to indicate where estimates are statistically reliable.

\subsection{Temperature Scaling}
\label{sec:methods:tempscaling}

Temperature scaling \citep{Guo2017} is a post-hoc calibration method 
that applies a single scalar parameter $T > 0$ to the pre-softmax 
logits of a classifier. For a classifier outputting probability vector 
$\mathbf{p}$, the calibrated probabilities are:

\begin{equation}
\hat{\mathbf{p}}_T = \mathrm{softmax}\!\left(\frac{\log \mathbf{p}}{T}\right)
\label{eq:tempscaling}
\end{equation}

Values of $T < 1$ increase confidence (sharpening the distribution); 
values of $T > 1$ decrease confidence (flattening toward uniform). 
The optimal $T$ is found by minimising ECE on the evaluation set 
using bounded scalar optimisation over $T \in [0.1, 10.0]$.

For Fink's binary classifiers, we apply an equivalent logit-space 
temperature scaling:

\begin{equation}
\hat{p}_T = \sigma\!\left(\frac{\mathrm{logit}(p)}{T}\right)
\label{eq:tempscaling_binary}
\end{equation}

where $\sigma$ is the sigmoid function and $\mathrm{logit}(p) = 
\log(p / (1-p))$.

Fitting $T$ on the same data used to evaluate it yields an optimistic
estimate. Our primary results therefore use held-out evaluation
throughout: 5-fold stratified cross-validation ($T$ fitted on four
folds by NLL minimisation, ECE evaluated on the held-out fold, pooled
across folds), corroborated by an independent 60/40
calibration/test split. We additionally report the same-set fit as an
explicit \emph{oracle upper bound}, labelled as such, for comparison
with the held-out numbers.

For the Fink classifiers we also fit isotonic regression
\citep{Zadrozny2002} under the same 5-fold cross-validation — a
substantially more flexible monotone calibrator than scalar
temperature — so that failure of temperature scaling alone is not
over-interpreted as failure of all post-hoc methods. Isotonic results
must be read jointly with the ranking metrics of
Section~\ref{sec:methods:ranking}: a calibrator applied to an
uninformative score can achieve low ECE by collapsing predictions
toward the base rate.

\subsection{Probability Renormalisation}
\label{sec:methods:renorm}

ALeRCE's classifier outputs probabilities across 15 classes, of which 
we evaluate 4 (excluding TDE, which is absent from ALeRCE's taxonomy). 
To assess whether apparent miscalibration arises from probability 
dilution across unevaluated classes, we compute a renormalised ECE 
by projecting the 15-class probability vector onto the 4 evaluated 
classes and renormalising to sum to unity. The difference between 
15-class and renormalised ECE quantifies the contribution of probability 
dilution to the observed miscalibration.